%!TeX program = xelatex
\PassOptionsToPackage{quiet}{fontspec}
\documentclass[10.5pt,compsoc,UTF8]{CjC}
\usepackage{thesis}
%===============================%信息设置
\mytitle{与茶共舞的未来：从小罐茶看茶产业的智能变革}
\myauthor{杨逍宇}
\mystudentid{3220105453}
\mycourse{《模电》}
%===============================%页眉设置
\headevenname{\mbox{\quad} \hfill  \mbox{\zihao{-5}{ \hfill \themycourse  } \hspace {50mm} \mbox{\today}}}%
\headoddname{\themyauthor \themystudentid \hfill \themycourse}%

\begin{document}
\setpage
\mytitlepage{\themytitle}{\themyauthor}{\themystudentid}
%=====================================================%摘要
\myabstract{
本文从现代工业应用的视角，探讨茶产业的未来发展。鉴于传统的茶叶生产方式在面临如需求变化、提升生产效率和市场占有率等方面的挑战时存在局限性，智能工厂的概念及其在茶产业应用被提出。本文试图通过对小罐茶智能工厂的剖析，探讨现代茶产业如何应对传统生产方式在面对市场变革和效率提升等挑战时的局限性。本文选取了知名的茶企小罐茶作为案例，介绍了其智能工厂在真实生产流程中如何具体运作，以及这种转型是如何解决茶产业的核心问题，有力地推动了中国茶叶的全球竞争力。
}
\vspace {5mm}
\keyword{智能工厂；小罐茶；茶工业}
%=====================================================%正文
\section{引言}



\section{第一节}
\section{第二节}
\cite{smith2019deep}
\newpage
\reference
\end{document}


